
\documentclass[uplatex,10pt,dvipdfmx,nomag]{jsarticle}

\usepackage[a4paper,top=20mm,bottom=24mm,left=22mm,right=22mm,headheight=14pt,headsep=6mm,footskip=10mm]{geometry}
\usepackage{amsmath,amssymb}
\usepackage{booktabs,tabularx,array,longtable,multicol}
\usepackage{enumitem}
\usepackage{fancyhdr}
\usepackage{lastpage}
\usepackage{titlesec}
\usepackage[dvipdfmx,bookmarks=true,bookmarksnumbered=true,hidelinks]{hyperref}
\usepackage{pxjahyper}

\setlength{\parindent}{1em}
\setlength{\parskip}{3pt}
\setlength{\tabcolsep}{5pt}
\renewcommand{\arraystretch}{1.18}
\linespread{1.05}
\raggedbottom
\sloppy
\emergencystretch=3em

\setlist[itemize]{leftmargin=1.8em,itemsep=.25ex,topsep=.35ex,parsep=0pt}
\setlist[enumerate]{leftmargin=2.0em,itemsep=.25ex,topsep=.35ex,parsep=0pt}

\titleformat{\section}{\normalfont\bfseries\large}{\thesection}{0.7em}{}
\titleformat{\subsection}{\normalfont\bfseries\normalsize}{\thesubsection}{0.6em}{}
\titleformat{\subsubsection}{\normalfont\bfseries\small}{\thesubsubsection}{0.5em}{}
\titlespacing*{\section}{0pt}{2.2ex plus .4ex}{.8ex}
\titlespacing*{\subsection}{0pt}{1.4ex plus .2ex}{.5ex}
\titlespacing*{\subsubsection}{0pt}{1.0ex plus .2ex}{.35ex}

\pagestyle{fancy}
\fancyhf{}
\fancyhead[L]{\footnotesize 第01章 ソフトウェア要求}
\fancyhead[R]{\footnotesize SWEBOK v4.0a 日本語まとめ}
\renewcommand{\headrulewidth}{0.25pt}
\fancyfoot[C]{\footnotesize \thepage/\pageref{LastPage}}

\newcommand{\term}[1]{\textbf{#1}}
\newcommand{\en}[1]{\textsf{#1}}
\newcommand{\Rule}{\par\smallskip\hrule\smallskip}
\newcommand{\MiniBox}[2]{%
  \par\smallskip\noindent\fbox{\begin{minipage}{0.94\linewidth}\textbf{#1}\par\smallskip #2\end{minipage}}\par\smallskip}
\newcommand{\TodoFigure}[2]{%
  \par\smallskip\noindent\fbox{\begin{minipage}{0.94\linewidth}
  \textbf{TODO（図）}\quad #1\par\smallskip
  \textbf{画像生成用プロンプト}\par
  {\footnotesize #2}
  \end{minipage}}\par\smallskip}
\newcommand{\SmallHead}[1]{\par\smallskip\noindent\textbf{#1}\quad}

\renewcommand{\contentsname}{目次}
\setcounter{tocdepth}{2}

\begin{document}

\thispagestyle{empty}
\begin{center}
{\Large\bfseries 第01章 ソフトウェア要求\par}
\vspace{1mm}
{\normalsize Software Requirements の日本語まとめ\par}
\vspace{1mm}
{\footnotesize SWEBOK Guide v4.0a Chapter 01 を初学者向けに再構成\par}
\vspace{2mm}
{\footnotesize 作成日：\today\par}
\end{center}
\Rule
\noindent\textbf{この資料の主張}

ソフトウェア要求は、単なる「欲しい機能の一覧」ではない。要求は、現実世界の問題を解くためにソフトウェアが満たすべき条件、能力、制約を言葉やモデルで表したものである。要求が曖昧なまま設計・実装・テストへ進むと、後で直す範囲が大きくなる。したがって、要求を「引き出す」「分析する」「仕様として書く」「妥当か確認する」「変化を管理する」という一連の活動として扱うことが重要である。

\MiniBox{この資料の読み方}{専門用語は原則として日本語で示し、必要に応じて英語を括弧内に併記する。英語名は、元資料やツール上の名称を確認したいときの手がかりとして残す。初学者が前提知識なしに読めるように、各節では「何のための概念か」「実務では何を見ればよいか」を重視する。}

\vspace{3mm}
\tableofcontents
\clearpage

\section*{全体概要：この章で学ぶこと}
\addcontentsline{toc}{section}{全体概要：この章で学ぶこと}

この章は、ソフトウェア開発の最初だけでなく、設計・実装・テスト・保守まで影響する「要求」を扱う。ここでいう要求とは、利用者や顧客の希望だけでなく、業務規則、法令、標準、運用上の制約、性能や安全性の水準、予算や納期の制約まで含み得る。

\SmallHead{到達目標}
この章を読み終えると、次のことができるようになる。
\begin{itemize}
\item 「要求」「機能要求」「非機能要求」「技術制約」「サービス品質制約」を区別できる。
\item 要求をどこから集めるか、どのように引き出すかを説明できる。
\item 要求が曖昧・不完全・矛盾している場合に、何を確認すべきかを説明できる。
\item 要求を文章・受け入れ基準・モデルとして記録する方法の違いを説明できる。
\item 要求の妥当性確認、変更管理、優先順位付け、追跡可能性の役割を説明できる。
\end{itemize}

\SmallHead{学習順序}
初学者は、まず「要求とは何か」と「要求の分類」を押さえる。その後、「要求をどう集めるか」「どう分析するか」「どう書くか」「どう確認するか」「どう変化に対応するか」という流れで読むとよい。

\TodoFigure{第01章の全体地図を入れる。}{白背景、モノクロ、A4本文幅に収まる横長の学習用図解を作成する。中央に「ソフトウェア要求」を置き、周囲に「基礎」「要求獲得」「要求分析」「要求仕様化」「要求妥当性確認」「要求管理」「実務上の考慮」「要求ツール」を配置する。矢印は「合意を作る活動」と「合意を維持する活動」の二系統で表す。ラベルは日本語、細線、講義ノート風、余白を広めにする。}

\subsection*{最初に必要な用語}
\addcontentsline{toc}{subsection}{最初に必要な用語}

\begin{tabularx}{0.97\linewidth}{@{}p{.25\linewidth}X@{}}
\toprule
用語 & 初学者向けの説明 \\
\midrule
\term{ソフトウェア} & プログラムだけでなく、データ、設定、操作手順、関連文書を含むことがある。 \\
\term{要求} & 何かを達成するために満たすべき条件・能力・制約。 \\
\term{利害関係者} & 開発結果に影響を与える、または影響を受ける人・組織。顧客、利用者、運用担当、規制当局など。 \\
\term{システム} & ソフトウェア、ハードウェア、人、手順、設備などが組み合わさって目的を達成するもの。 \\
\term{仕様} & 要求や設計などを、関係者が共有できるように記録したもの。 \\
\term{妥当性確認} & 「これを作れば本当に相手の問題を解決できるか」を確認すること。 \\
\term{追跡可能性} & 要求が設計、コード、テスト、利用者文書などとどうつながるかをたどれること。 \\
\bottomrule
\end{tabularx}

\section*{導入：要求が重要な理由}
\addcontentsline{toc}{section}{導入：要求が重要な理由}

ソフトウェア要求は、二つの視点から理解する。第一に、要求は「ソフトウェア製品やプロジェクトに対するニーズと制約の表現」である。第二に、要求は「そのニーズと制約を作り、保ち、変化に対応する活動」である。

要求を誤ると、プロジェクトの費用、納期、品質に大きく影響する。ひとつの要求は複数の設計判断を生み、ひとつの設計判断は複数のコード判断を生み、さらに複数のテスト判断につながる。したがって、要求の欠落や誤解が後工程で見つかると、修正範囲が連鎖的に広がる。

現実のプロジェクトで特に多い問題は、\term{不完全性}と\term{曖昧性}である。不完全性とは、必要な要求や詳細が明らかになっていない状態である。曖昧性とは、同じ要求文を複数の意味に解釈できる状態である。どちらも、作ったものが「本来ほしかったもの」とずれる原因になる。

要求文書は、初期開発のためだけではない。保守担当者が後からコードを見ると、コードが「何をしているか」は実行や解析でわかる場合がある。しかし、コードが「何を意図しているか」は要求がなければわかりにくい。欠陥とは、多くの場合、ソフトウェアが意図されたことと実際にしていることの差である。要求は、その意図を長期にわたって伝える役割を持つ。

\MiniBox{重要}{要求活動は、開発の最初に一回だけ行う作業ではない。要求は、プロジェクト開始時に始まり、開発中に洗練され、運用・保守の期間にも参照・変更される。}

\section{ソフトウェア要求の基礎}

\subsection{ソフトウェア要求の定義}

\term{ソフトウェア要求}とは、利用者が問題を解決したり目的を達成したりするために必要な条件または能力である。また、契約、標準、仕様、法令などを満たすために、システムや構成要素が備えるべき条件または能力でもある。さらに、それらを文書やモデルとして表したものも要求と呼ぶ。

より実務的に言えば、ソフトウェア要求とは「何を満たせば、現実世界の問題を解決したと言えるか」を表す合意である。ここでいう「問題」は、業務の自動化、既存ソフトウェアの欠点の修正、機器の制御、利用者体験の改善、規制への対応など、さまざまである。

\MiniBox{定義}{ソフトウェア要求とは、ソフトウェアまたはソフトウェアを作るプロジェクトが満たすべき条件・能力・制約を、関係者が共有できるように表したものである。}

要求は通常、顧客、利用者、業務担当者から出る。しかし、それだけではない。規制当局、標準、既存システム、運用部門、サポート部門、開発組織、プロジェクト計画そのものも要求の源泉になる。

\subsection{ソフトウェア要求の分類}

要求は、すべてをひとまとめにすると扱いにくい。SWEBOKでは、ソフトウェア要求を大きく次のように整理する。

\begin{center}
\footnotesize
\begin{tabular}{c}
ソフトウェア要求 \\
$\downarrow$ \\
ソフトウェア製品要求 \quad / \quad ソフトウェアプロジェクト要求 \\
$\downarrow$ \\
機能要求 \quad / \quad 非機能要求 \\
$\downarrow$ \\
技術制約 \quad / \quad サービス品質制約
\end{tabular}
\end{center}

\TodoFigure{要求分類の階層図を入れる。}{白背景、モノクロ、A4本文幅に収まる縦長の階層図を作成する。最上位に「ソフトウェア要求」。第2層に「ソフトウェア製品要求」と「ソフトウェアプロジェクト要求」。製品要求から「機能要求」と「非機能要求」へ分岐。非機能要求から「技術制約」と「サービス品質制約」へ分岐。各箱に短い説明を添える。講義ノート風、細線、印刷で読みやすい図。}

分類の目的は、単に用語を覚えることではない。分類すると、要求の源泉、引き出し方、分析方法、仕様化の方法、確認すべき相手が見えやすくなる。

\subsection{ソフトウェア製品要求とソフトウェアプロジェクト要求}

\term{ソフトウェア製品要求}は、完成するソフトウェアそのものが満たすべき性質である。画面、処理、データ、外部連携、性能、信頼性、セキュリティなどが含まれる。

\term{ソフトウェアプロジェクト要求}は、ソフトウェアを作るプロジェクトの進め方を制約する要求である。費用、納期、人員、テスト環境、データ移行、利用者教育、保守準備などが例である。プロジェクト要求は、プロジェクト憲章、契約、計画書などに記録されることが多い。

\begin{tabularx}{0.97\linewidth}{@{}p{.28\linewidth}X@{}}
\toprule
区分 & 見分け方 \\
\midrule
\term{製品要求} & 完成したソフトウェアがどう振る舞うか、どの品質を満たすかを述べる。 \\
\term{プロジェクト要求} & 開発の費用、納期、人員、進め方、移行、教育などを縛る。 \\
\bottomrule
\end{tabularx}

\MiniBox{例}{「利用者は注文履歴を検索できる」は製品要求である。「本番移行は年度末までに完了する」はプロジェクト要求である。}

\subsection{機能要求}

\term{機能要求}とは、ソフトウェアが提供すべき観測可能な振る舞いを表す要求である。言い換えると、「ソフトウェアが何をするか」を述べる要求である。

機能要求は、二つの観点で理解しやすい。第一に、\term{方針}である。方針とは、常に守られるべき規則である。たとえば銀行システムでは、「口座は少なくとも一人の所有者を持つ」「口座残高は負にならない」といった要求が該当する。第二に、\term{プロセス}である。プロセスとは、実行される業務手続きである。入金、出金、振込、予約、取消、検索、承認などが該当する。

\MiniBox{定義}{機能要求とは、ソフトウェアが実行すべき処理、守るべき業務規則、利用者に提供すべき振る舞いを述べる要求である。}

\subsection{非機能要求}

\term{非機能要求}とは、ソフトウェアの実装技術や運用品質を制約する要求である。「何をするか」ではなく、「どのような条件で、どの水準で、どの技術的制約のもとで実現するか」に関わる。

非機能要求には、性能、応答時間、信頼性、正確性、可用性、拡張性、保存容量、プラットフォーム、利用するデータベースなどが含まれる。非機能要求はさらに、\term{技術制約}と\term{サービス品質制約}に分けると理解しやすい。

\MiniBox{注意}{「非機能要求」は「重要でない要求」という意味ではない。むしろ、性能、信頼性、セキュリティ、安全性など、失敗したときの影響が大きい要求が多い。}

\subsection{技術制約}

\term{技術制約}とは、特定の自動化技術やインフラを使うこと、または使わないことを指定する要求である。ここでは「名前のある技術を指定しているか」に注目すると見分けやすい。

\begin{tabularx}{0.97\linewidth}{@{}p{.25\linewidth}X@{}}
\toprule
対象 & 例 \\
\midrule
プラットフォーム & Windows、macOS、Android、iOS など。 \\
プログラミング言語 & Java、C++、C\#、Python など。 \\
ブラウザ & Chrome、Safari、Edge など。 \\
データベース & Oracle、SQL Server、MySQL など。 \\
禁止事項 & ポインタ利用禁止、特定暗号方式の禁止、外部クラウド利用禁止など。 \\
\bottomrule
\end{tabularx}

技術制約は、既存システムとの連携、組織標準、保守要員のスキル、セキュリティ方針、ライセンス方針などから生じることが多い。

\subsection{サービス品質制約}

\term{サービス品質制約}とは、特定の技術名を指定せず、ソフトウェアが達成すべき品質水準を表す要求である。サービス品質という表現は、利用者や運用者が受け取るサービスの品質として考えるとわかりやすい。

\begin{tabularx}{0.97\linewidth}{@{}p{.26\linewidth}X@{}}
\toprule
品質水準 & 要求文の例 \\
\midrule
応答時間 & 通常時、検索結果を2秒以内に返す。 \\
スループット & 1分あたり1万件のイベントを処理する。 \\
正確性 & 金額計算は定められた丸め規則に従う。 \\
信頼性 & 月間稼働率99.9\%以上を維持する。 \\
拡張性 & 利用者数が10倍になっても構成変更で対応できる。 \\
セキュリティ & 権限のない利用者は個人情報を閲覧できない。 \\
\bottomrule
\end{tabularx}

\subsection{なぜこのように分類するのか}

要求を分類する理由は、実務上は非常に明確である。

\begin{itemize}
\item 要求の種類によって、要求の源泉が異なる。
\item 源泉が異なれば、聞き出し方も異なる。
\item 種類によって、分析方法、仕様化方法、確認すべき相手が異なる。
\item 要求の種類によって、設計・実装・テストへの影響が異なる。
\end{itemize}

さらに、分類は複雑さを下げる。業務規則を議論している最中に、同時にデータベース製品や処理速度を議論すると、論点が混ざる。機能要求、技術制約、サービス品質制約を分ければ、業務の専門家は業務の観点に集中でき、技術者は技術の観点に集中できる。

\MiniBox{完全技術フィルタ}{機能要求と非機能要求の境界が曖昧なときは、「無限に速く、記憶容量が無限で、費用がゼロで、故障しない理想的なコンピュータがあっても、その要求を書く必要があるか」と問う。必要なら機能要求である。不要なら、技術や品質に関する制約、つまり非機能要求である。}

\TodoFigure{完全技術フィルタの判定フローを入れる。}{白背景、モノクロ、A4本文幅のフローチャートを作成する。開始は「候補要求」。判定ひし形に「理想的なコンピュータでも必要か？」と書く。Yes は「機能要求」。No は「非機能要求」。さらに「特定技術名を指定するか？」で分岐し、Yes は「技術制約」、No は「サービス品質制約」。日本語ラベル、細線、講義ノート風。}

\subsection{システム要求とソフトウェア要求}

\term{システム要求}は、ソフトウェアだけでなく、ハードウェア、ファームウェア、人、情報、手順、設備、サービス、支援要素を含むシステム全体に対する要求である。自動運転車、医療機器、工場制御、航空機システムのように、複数の要素が組み合わさる場合には、システム要求とソフトウェア要求を分けて扱うことが重要である。

\term{ソフトウェア要求}は、その大きなシステムの中で、ソフトウェア要素が満たすべき要求である。たとえば「車両は障害物を検知して安全に停止する」はシステム要求であり、その一部として「認識ソフトウェアは前方物体を一定時間内に分類する」といったソフトウェア要求が割り当てられる。

\TodoFigure{システム要求からソフトウェア要求への割当図を入れる。}{白背景、モノクロ、A4本文幅の概念図を作成する。大きな外枠を「システム」とし、内側に「ハードウェア」「ソフトウェア」「人」「手順」「設備」の箱を配置する。上から「システム要求」が外枠全体に向かう矢印を描き、その一部が「ソフトウェア要求」としてソフトウェア箱へ割り当てられる矢印を描く。日本語ラベル、講義ノート風。}

\subsection{派生要求}

\term{派生要求}とは、外部の利害関係者から直接出された要求ではなく、開発チーム内部の判断から生まれ、下位チームや構成要素に課される要求である。

たとえば、アーキテクトが「パイプ・アンド・フィルタ構造を採用する」と決めたとする。この決定は、外部顧客から見ると設計判断である。しかし、個々のフィルタを作るサブチームから見ると、「この形式で入力を受け取り、この形式で出力する」という要求になる。このように、要求か設計かは視点に依存する。

\MiniBox{視点依存性}{同じ記述でも、上位の視点では設計判断、下位の視点では要求になることがある。要求とは「その作業を始める人にとって、満たさなければならない入力条件」として理解するとよい。}

\subsection{ソフトウェア要求活動}

要求活動は、大きく\term{要求開発}と\term{要求管理}に分けられる。

\term{要求開発}は、「何を作るか」について合意を作る活動である。具体的には、\term{要求獲得}、\term{要求分析}、\term{要求仕様化}、\term{要求妥当性確認}を含む。

\term{要求管理}は、その合意を時間とともに維持する活動である。要求は開発中にも保守中にも変化するため、不要な要求の整理、変更要求の扱い、スコープと予算・納期の調整が必要になる。

\TodoFigure{要求開発と要求管理の活動図を入れる。}{白背景、モノクロ、A4本文幅の工程図を作成する。上段に「要求開発」と書き、その下に「獲得」「分析」「仕様化」「妥当性確認」の箱を左から右へ並べる。下段に「要求管理」と書き、「要求整理」「変更管理」「範囲調整」の箱を並べる。要求開発は「合意を作る」、要求管理は「合意を維持する」と注記する。各活動の間に反復を示す矢印を入れる。}

\section{要求獲得}

\term{要求獲得}とは、候補となる要求を表面化させる活動である。英語では \,\en{Requirements Elicitation}\, と呼ばれる。直訳すれば「引き出し」であり、相手が最初から明確に言語化できるとは限らないニーズを、質問、観察、試作、資料調査などで明らかにする活動である。

要求獲得が不十分だと、不完全な要求が残る。完全性を保証することは難しいが、適切な源泉を見つけ、複数の技法を組み合わせることで、不完全性を減らせる。

\subsection{要求の源泉}

要求は多くの源泉から来る。人に由来する源泉と、人以外に由来する源泉の両方を考える必要がある。

\begin{tabularx}{0.97\linewidth}{@{}p{.27\linewidth}X@{}}
\toprule
源泉 & 例 \\
\midrule
\term{顧客} & ソフトウェア開発に費用を出す人・組織。 \\
\term{利用者} & ソフトウェアを直接または間接に使う人。利用頻度、権限、熟練度で分類できる。 \\
\term{業務専門家} & 対象業務、規制、現場作業に詳しい人。 \\
\term{運用担当} & 本番環境、監視、障害対応、バックアップなどに責任を持つ人。 \\
\term{サポート担当} & 利用者からの問い合わせや障害報告を受ける人。 \\
\term{規制当局・専門団体} & 法令、標準、ガイドラインを通じて要求を課す組織。 \\
\term{開発者} & 技術的制約、保守性、開発効率に関する要求を示すことがある。 \\
\bottomrule
\end{tabularx}

また、要求は人からだけではなく、文書、既存システム、業務方針、計算環境、競合製品、障害履歴、標準、法規制からも得られる。

\MiniBox{利害関係者分析}{利害関係者分析とは、関係する人や組織をできるだけ漏れなく見つけ、代表性の偏りを減らす活動である。声の大きい利用者だけから要求を集めると、少数派や運用担当の重要な要求を見落としやすい。}

\subsection{要求獲得技法}

要求獲得には多くの技法がある。技法は相手や目的によって使い分ける。

\begin{tabularx}{0.97\linewidth}{@{}p{.25\linewidth}X@{}}
\toprule
技法 & 何に向いているか \\
\midrule
\term{面接} & 個別の経験、例外処理、暗黙知を聞き出す。 \\
\term{会議・討議} & 複数の関係者で認識を合わせる。 \\
\term{促進付きワークショップ} & 利害関係者を集め、短時間で合意形成や要求整理を行う。 \\
\term{質問票・市場調査} & 多数の人から傾向を集める。 \\
\term{観察} & 現場作業や実際の利用状況を確認する。 \\
\term{徒弟型学習} & ソフトウェア技術者が実際の作業を体験し、暗黙知を得る。 \\
\term{試作} & 画面や操作感を具体化し、利用者から反応を得る。 \\
\term{ユーザーストーリーマッピング} & 利用者の行動の流れに沿って機能を整理する。 \\
\term{文献・標準調査} & 法規制、標準、既存仕様に基づく要求を見つける。 \\
\bottomrule
\end{tabularx}

\TodoFigure{要求の源泉と獲得技法の対応図を入れる。}{白背景、モノクロ、A4本文幅のマトリクス図を作成する。左列に「利用者」「業務専門家」「運用担当」「既存文書」「既存システム」「規制・標準」。上段に「面接」「観察」「ワークショップ」「試作」「文書調査」。交点に丸印を置き、どの源泉にどの技法が向くかを示す。日本語、講義ノート風、印刷で読みやすい。}

要求獲得は受け身の活動ではない。利用者は自分の作業をうまく説明できないことがある。重要な情報を当然のこととして言わないこともある。したがって、技術者は「言われたことを書くだけ」ではなく、質問し、観察し、仮説を作り、確認する必要がある。

\section{要求分析}

\term{要求分析}とは、獲得された候補要求の意味、妥当性、矛盾、実現可能性、影響を調べる活動である。要求は、最初から完成形で出てくるとは限らない。多くの場合、曖昧、過剰、矛盾、不足、または解決策の断片として現れる。分析の目的は、それを真の要求へ近づけることである。

\subsection{要求の性質}

よい要求には、個々の要求として望ましい性質と、要求集合として望ましい性質がある。

\subsubsection{個々の要求の望ましい性質}

\begin{tabularx}{0.97\linewidth}{@{}p{.26\linewidth}X@{}}
\toprule
性質 & 説明 \\
\midrule
\term{一意に解釈できる} & 読む人によって意味が変わらない。 \\
\term{テスト可能である} & 満たしたかどうかを確認できる。必要なら数値で示す。 \\
\term{拘束力がある} & 顧客がそれに価値を認め、満たされないことを受け入れない。 \\
\term{原子的である} & 一つの要求が一つの判断だけを表す。 \\
\term{真のニーズを表す} & 単なる思いつきや特定解ではなく、実際の問題に基づく。 \\
\term{利害関係者の語彙を使う} & 利用者や業務専門家が理解できる言葉で書く。 \\
\term{関係者に受け入れられる} & 重要な利害関係者が納得できる。 \\
\bottomrule
\end{tabularx}

\subsubsection{要求集合の望ましい性質}

\begin{tabularx}{0.97\linewidth}{@{}p{.26\linewidth}X@{}}
\toprule
性質 & 説明 \\
\midrule
\term{完全性} & 通常ケースだけでなく、境界条件、例外条件、セキュリティ上の必要事項も扱う。 \\
\term{簡潔性} & 不要な記述や重複が少ない。 \\
\term{内部整合性} & 要求同士が矛盾しない。 \\
\term{外部整合性} & 契約、法令、標準、業務文書などの外部資料と矛盾しない。 \\
\term{実現可能性} & 費用、納期、人員、技術制約の範囲で実現できる。 \\
\bottomrule
\end{tabularx}

\MiniBox{5つのなぜ}{要求が解決策の形で出てきたときは、「なぜそれが必要なのか」を繰り返し問う。目的は、相手が本当に解決したい問題に到達することである。}

\subsection{サービス品質要求の経済性}

サービス品質制約は、単に「高ければよい」ものではない。たとえば、応答時間を極限まで短くする、稼働率を極限まで高くする、処理能力を極限まで大きくするには、費用が増える。したがって、要求は価値と費用の関係で考える必要がある。

\term{失敗点}は、それを下回るとシステムとして価値がほとんどなくなる水準である。\term{完全点}は、それを超えても利用者に追加価値がほとんど出ない水準である。現実的には、要求値はこの間で決まる。最もよい水準は、価値と費用の差が最大になる点である。

\TodoFigure{サービス品質制約の価値・費用曲線を入れる。}{白背景、モノクロ、A4本文幅に収まる2段の概念図を作成する。上段は横軸「性能水準」、縦軸「価値」。右上がりの曲線を描き、左側に「失敗点」、右側に「完全点」を示す。下段は横軸「性能水準」、縦軸「金額」。価値曲線と階段状の費用線を描き、差が最大になる点を「最も費用対効果が高い水準」と示す。日本語ラベル、講義ノート風。}

サービス品質制約同士には、よい関係と悪い関係がある。たとえば、コードを単純にすると信頼性と変更容易性が同時に上がることがある。一方、実行速度を極端に高める最適化が、コードを複雑にして変更容易性を下げることもある。

\subsection{形式的分析}

\term{形式的分析}とは、数学的に意味が定義された言語やモデルを使って、要求を精密に表し、性質を検査する方法である。たとえば、デッドロックが起きないこと、特定の状態へ必ず到達すること、矛盾した状態が存在しないことなどを確認するために使われる。

形式的分析の利点は、曖昧さを減らせること、要求を論理的に推論できること、通常のレビューでは見落としやすい性質を確認できることである。一方で、記法を学ぶ負担があり、すべてのプロジェクトに適するわけではない。安全性や信頼性が非常に重要な領域で特に価値が高い。

\subsection{要求の衝突への対応}

利害関係者が多いほど、要求の衝突は起こりやすい。たとえば、利用者は高機能を求め、運用担当は単純で安定した運用を求め、経営側は費用や納期を抑えたいと考えることがある。

衝突への対応には、主に二つの方向がある。第一に、関係者間で交渉し、合意を形成することである。技術者が独断で決めると、責任や契約上の問題が残ることがある。第二に、\term{製品群開発}の考え方を使い、すべての利用者に共通する不変要求と、利用者や市場によって変わる可変要求を分けることである。

\MiniBox{例}{天気アプリで、ある利用者は摂氏表示を求め、別の利用者は華氏表示を求める。この場合、温度を表示することは共通要求であり、単位は可変要求である。}

\section{要求仕様化}

\term{要求仕様化}とは、要求を記録し、記憶し、伝達できる形にする活動である。仕様化の方法は一つではない。自然言語で書く、構造化したテンプレートで書く、受け入れ基準として書く、モデルで表すなど、複数の方法がある。

どの方法を選ぶかは、プロジェクトの性質によって変わる。業務領域への理解、リスクの大きさ、チームの地理的分散、外注の有無、契約上の要求、要求に基づくテストの必要性、保守期間中の担当者交代などを考慮する必要がある。

\MiniBox{基本方針}{要求仕様は、読む人を意識して作る。利用者、顧客、設計者、開発者、テスト担当、保守担当は、それぞれ必要とする情報が違う。}

\subsection{非構造化自然言語による要求仕様化}

\term{非構造化自然言語}とは、通常の文章で要求を書く方法である。たとえば「システムは、利用者が注文履歴を検索できるようにする」といった文である。

利点は、誰でも読み書きしやすいことである。欠点は、曖昧さや抜け漏れが発生しやすいことである。自然言語は柔軟であるため、同じ文が複数の意味に読めることがある。

\subsection{構造化自然言語による要求仕様化}

\term{構造化自然言語}とは、自然言語を使いながら、書き方に制約を加える方法である。代表例は、\term{主体・動作形式}である。「いつ」「誰が」「何をする」「どの条件で」という形にそろえると、曖昧さを減らせる。

\MiniBox{例}{「注文が出荷されたとき、システムは請求書を作成する。ただし、注文条件が前払いの場合を除く。」この文では、きっかけ、主体、動作、例外条件が明示されている。}

他にも、\term{ユースケース仕様}、\term{ユーザーストーリー}、\term{決定表}などが構造化自然言語の例である。ユーザーストーリーは、「ある役割として、ある能力がほしい。なぜなら、ある便益を得たいからである」という形式で、利用者視点の価値を表す。

\subsection{受け入れ基準に基づく要求仕様化}

\term{受け入れ基準}とは、要求が満たされたと判断するための条件である。要求を受け入れ基準の形で書くと、テストと要求が強く結びつく。

\term{受け入れテスト駆動開発}（ATDD）は、機能を作る前に、その機能が完成したと判断するテストを合意する方法である。\term{振る舞い駆動開発}（BDD）は、利用者の振る舞いを「前提、事象、結果」の形で書く方法である。

\MiniBox{BDDの例}{前提：口座残高が500円で、カードが有効である。事象：利用者が100円の出金を要求する。結果：ATMは100円を払い出し、口座残高は400円になる。}

この方法は、曖昧さの低減に有効である。一方で、どのシナリオを書くべきかを見落とすと不完全性は残る。そのため、境界値分析、組合せテスト、同値分割などのテスト技法と組み合わせるとよい。

\subsection{モデルに基づく要求仕様化}

\term{モデル}とは、複雑な対象を理解しやすくするために、必要な部分だけを抽象化して表したものです。要求仕様化では、文章だけでなく、図や形式的な記法を使って要求を表すことがある。UMLやSysMLの一部は、そのために使われる。

モデルは大きく二種類に分けられる。\term{構造モデル}は、概念、データ、関係、業務規則を表す。\term{振る舞いモデル}は、処理の流れ、利用者とのやり取り、状態遷移、イベントへの反応を表す。

\begin{tabularx}{0.97\linewidth}{@{}p{.24\linewidth}X@{}}
\toprule
モデルの種類 & 例 \\
\midrule
構造モデル & クラス図、概念データモデル、実体関連図。 \\
振る舞いモデル & ユースケース図、シーケンス図、状態遷移図、活動図、データフロー図。 \\
\bottomrule
\end{tabularx}

モデルには、手描きのラフな図から、数学的意味を厳密に持つ形式モデルまで幅がある。形式性が高いほど曖昧さは減るが、作成と理解の負担は増える。

\TodoFigure{要求仕様化技法の比較図を入れる。}{白背景、モノクロ、A4本文幅の比較図を作成する。横軸に「曖昧さが残りやすい」から「精密」の尺度、縦軸に「作成が容易」から「作成が難しい」の尺度を置く。点として「非構造化自然言語」「構造化自然言語」「受け入れ基準」「モデル」「形式モデル」を配置する。日本語ラベル、講義ノート風。}

\subsection{要求の追加属性}

要求文そのものに加えて、要求を管理するための属性を付けると有用である。

\begin{tabularx}{0.97\linewidth}{@{}p{.25\linewidth}X@{}}
\toprule
属性 & 目的 \\
\midrule
識別子 & 要求を一意に参照し、追跡する。 \\
説明 & 要求文だけでは足りない背景を補足する。 \\
根拠 & なぜその要求が必要かを説明する。 \\
源泉 & 誰または何から出た要求かを示す。 \\
種類 & 機能要求、技術制約、サービス品質制約などを示す。 \\
依存関係 & 他の要求との前後関係や依存を示す。 \\
衝突 & 他の要求と矛盾する可能性を示す。 \\
受け入れ基準 & 完了判定に使う条件を示す。 \\
優先度 & 重要度や実施順序を示す。 \\
安定性 & 変更される可能性を示す。 \\
変更履歴 & いつ、誰が、なぜ変更したかを残す。 \\
\bottomrule
\end{tabularx}

\subsection{増分的仕様化と包括的仕様化}

要求を明示的に文書化するプロジェクトでは、\term{増分的仕様化}と\term{包括的仕様化}の二つの考え方がある。

増分的仕様化では、今回の版で追加・変更・削除された差分だけを記録する。文書量は小さくなりやすいが、全体像を理解するには過去の差分を追う必要がある。

包括的仕様化では、各版にすべての要求を含める。文書量は増えるが、その版だけで全体を理解できる。実務では、マイナー版は差分、メジャー版は包括的に書くという組合せもある。

\section{要求妥当性確認}

\term{要求妥当性確認}とは、現在理解されている要求が、利害関係者の真のニーズを表しているかを確認する活動である。英語では \,\en{Requirements Validation}\, と呼ばれる。

ここで大切なのは、「正しく書かれているか」だけではなく、「そもそもこれを作れば問題が解決するか」を見ることである。要求文が文法的に正しくても、利用者の実際の困りごとを解決しないなら妥当ではない。

\MiniBox{妥当性確認で問うこと}{要求はこの時点で必要なものを十分に含んでいるか。不要な要求はないか。理解しやすいか。矛盾していないか。標準や契約に合っているか。}

\subsection{要求レビュー}

\term{要求レビュー}は、要求文書やモデルを人が読んで確認する方法である。もっとも一般的な妥当性確認手段である。

レビューには複数の視点が必要である。顧客や利用者は、自分たちのニーズが正しく表されているかを見る。要求の専門家は、文書が明確で、標準に沿っているかを見る。設計者や開発者は、設計・実装に必要な情報があるかを見る。テスト担当は、テスト可能かを見る。

\subsection{シミュレーションと実行}

形式的またはモデルベースの要求仕様では、仕様を手作業またはツールで\term{実行}または\term{シミュレーション}できることがある。利用者は、文章を読むよりも、シナリオに沿って動く様子を見る方が理解しやすい場合がある。

\subsection{プロトタイピング}

\term{プロトタイプ}とは、本格的な実装の前に作る試作品である。画面の動き、操作感、処理の流れ、重要な技術的可能性などを具体的に示すために使う。

プロトタイプは、技術者の解釈を利害関係者に見せ、誤解を早く発見するのに有効である。ただし、見た目の細部や一時的な品質問題に関心が向きすぎると、本来確認したい要求から注意がそれることがある。

\TodoFigure{要求妥当性確認の三手法を比較する図を入れる。}{白背景、モノクロ、A4本文幅の比較表風図を作成する。列は「レビュー」「シミュレーション・実行」「プロトタイプ」。行は「向いている対象」「強み」「注意点」。各セルは短い日本語。講義ノート風、細線、印刷で読みやすい。}

\section{要求管理活動}

要求開発が「合意を作る活動」だとすれば、\term{要求管理}は「合意を維持する活動」である。要求は開発中にも保守中にも変わる。したがって、不要な要求を減らし、変更を統制し、範囲と制約を整合させる必要がある。

\subsection{要求整理・削減}

\term{要求整理・削減}は、英語の \,\en{Requirements Scrubbing}\, に対応する。目的は、利害関係者のニーズを満たすために必要な、できるだけ小さく、単純な要求集合を得ることである。

削減の対象は、範囲外の要求、費用対効果が低い要求、重要度が低い要求、過度に複雑な要求である。要求が減れば、設計、実装、テスト、保守の負担も減る。

\subsection{要求変更管理}

\term{要求変更管理}とは、合意済みの要求を変更する手順を管理する活動である。計画駆動型の開発では、変更要求、影響分析、承認・却下・保留の判断、関係者への通知、完了までの追跡が明示的に行われることが多い。

アジャイル型の開発では、変更要求がプロダクトバックログの項目として扱われ、優先度が高くなったときに反復に取り込まれることが多い。どちらの場合も、変更を受け入れるとは、納期、費用、人員、または他の範囲への影響を受け入れることを意味する。

\subsection{範囲調整}

\term{範囲調整}とは、作るべき要求の範囲が、費用、納期、人員などの制約を超えないように調整する活動である。範囲が大きすぎる場合には、要求を減らす、予算や人員を増やす、納期を延ばす、またはそれらを組み合わせる必要がある。

可能であれば、範囲調整は感覚ではなく、機能規模、工数、リスクなどの測定値に基づいて行うとよい。

\TodoFigure{要求変更管理の流れを入れる。}{白背景、モノクロ、A4本文幅のフローチャートを作成する。「変更要求」から始まり、「影響分析」「判断」「通知」「実施」「完了確認」へ進む。判断から「承認」「却下」「保留」の三分岐を描く。各箱に短い説明を入れる。日本語ラベル、講義ノート風。}

\section{実務上の考慮}

\subsection{要求プロセスの反復性}

現実の要求活動は、一回で終わらない。最初は広く浅く要求を集め、次に重要な部分を深く分析し、その後また不足が見つかって追加で獲得することが多い。要求は幅と深さの両方を持つため、反復的に進めるのが自然である。

ウォーターフォール型では、要求フェーズで多くの要求を先に固める。反復型やアジャイル型では、初期に高レベル要求を押さえ、その後、開発する順に詳細化する。重要なのは、ライフサイクルの名前ではなく、対象ソフトウェアの性質に合う要求活動を選ぶことである。

\subsection{要求の優先順位付け}

\term{要求の優先順位付け}は、限られた費用と時間の中で、価値の高い要求から扱うために必要である。優先度は、価値だけでなく、未実現だった場合の不満、実装費用、保守費用、技術リスク、利用されないリスクなどを考慮して決める。

\MiniBox{Kanoモデルの考え方}{ある機能があると満足度が上がる場合と、ないと強い不満が出る場合は異なる。たとえばメールソフトでは、迷惑メールフィルタはあると嬉しい機能かもしれないが、添付ファイルを扱えないと利用者は強く不満を持つ。優先度は満足と不満の両方から考える。}

優先度の表し方には、「必須」「望ましい」「あればよい」のような段階、1から10の数値、優先順位リストなどがある。実務では、細かな数値差よりも、同程度の優先度を持つ要求群を見つけることが重要である。

\subsection{要求の追跡可能性}

\term{追跡可能性}とは、要求がどこから来て、どの設計、コード、テスト、文書につながるかをたどれることです。これは二つの目的で役立つ。

第一に、成果物間の整合性確認である。ある要求に対応する設計要素がなければ、その要求は満たされていない可能性がある。逆に、ある設計要素に対応する要求がなければ、その設計要素は不要か、要求が不足している可能性がある。

第二に、変更の影響分析である。ある要求が変わったとき、関連するソフトウェア要求、設計要素、コード、テストケース、利用者文書をたどることで、変更に必要な作業範囲を見積もりやすくなる。

\TodoFigure{要求追跡の連鎖図を入れる。}{白背景、モノクロ、A4本文幅の横長図を作成する。左から「要求源泉」「ソフトウェア要求」「設計」「コード」「テストケース」「利用者文書」を箱で並べ、矢印でつなぐ。上に「前方追跡」、下に「後方追跡」の矢印を描く。変更時には影響範囲をたどることを点線で示す。日本語ラベル、講義ノート風。}

\subsection{要求の安定性と変動性}

\term{要求の安定性}とは、要求が変わりにくい度合いである。\term{変動性}とは、要求が変わりやすい度合いである。

銀行アプリケーションで「口座に利息を計算して付与する」は比較的安定した要求である。一方、税制に基づく特定口座の扱いは、法改正で変わる可能性がある。変わりやすい要求を早く見つけると、変更に強い設計を選びやすくなる。

\subsection{要求の測定}

要求を測定する理由は、規模、費用、進捗、品質を管理するためである。\term{機能規模測定}は、機能要求の大きさを測る技法である。アジャイル開発で使われるストーリーポイントも、要求の大きさを表す測度として使われることがある。

要求品質の指標は、曖昧さ、テスト可能性、重複、矛盾、欠落、変更頻度などから作れる。これらを測ることで、要求プロセスの弱点を改善しやすくなる。

\subsection{要求プロセスの品質と改善}

要求プロセスの品質は、ソフトウェア製品の費用、納期、顧客満足度に大きく影響する。要求プロセスを改善するには、プロセス改善モデルや品質標準との対応、測定、ベンチマーク、改善計画、実施、振り返りが必要である。

セキュリティの観点では、機密性、完全性、可用性をどのように要求へ反映し、改善するかも重要である。

\section{ソフトウェア要求ツール}

要求活動は、紙や表計算ソフトだけでも始められる。しかし、要求が多くなり、変更や追跡が必要になると、専用ツールの支援が重要になる。

\subsection{要求管理ツール}

\term{要求管理ツール}は、要求の属性管理、追跡可能性、文書生成、変更管理などを支援する。要求が増えるほど、手作業での追跡は難しくなる。特に、規制対応や安全性が重要な領域では、ツールの価値が高い。

\subsection{要求モデリングツール}

\term{要求モデリングツール}は、要求モデルを作成、変更、公開するためのツールである。図を描くだけでなく、構文チェック、整合性確認、完全性確認、シミュレーションなどを支援するものもある。

形式的分析では、定理証明器やモデル検査器のようなツールが使われる。ただし、ツールを使えば自動的にすべてが証明できるわけではなく、形式的な推論に関する知識が必要になる。

\subsection{機能テストケース生成ツール}

要求仕様の形式性が高いほど、機能テストケースを機械的に導きやすくなる。たとえば、BDDのシナリオはテストケースへ変換しやすい。状態モデルでは、定義された遷移から正のテストケースを作り、存在しない状態とイベントの組合せから負のテストケースを作れる。

ただし、一般にはツールが入力データを生成できても、期待結果の判断には業務知識が必要になることがある。

\TodoFigure{要求ツールの分類図を入れる。}{白背景、モノクロ、A4本文幅の三分割図を作成する。左から「要求管理ツール」「要求モデリングツール」「機能テストケース生成ツール」。各箱の下に代表的な支援内容を3つずつ書く。管理ツールは「属性」「追跡」「変更管理」、モデリングツールは「図化」「整合性確認」「シミュレーション」、テスト生成ツールは「入力生成」「シナリオ変換」「期待結果確認支援」とする。日本語、講義ノート風。}

\section{トピックと参考文献の対応}

この表は、SWEBOKが示す「トピック対参考資料の行列」を、学習用に簡略化したものである。詳しく学ぶときは、各トピックに対応する文献を読む。

\small
\begin{longtable}{@{}p{.38\linewidth}p{.55\linewidth}@{}}
\toprule
トピック & 主な参考資料 \\
\midrule
\endhead
1. ソフトウェア要求の基礎 & Wiegers and Beatty, \textit{Software Requirements}; Sommerville, \textit{Software Engineering} \\
1.1 要求の定義 & Wiegers and Beatty; Sommerville; ISO/IEC/IEEE 24765 \\
1.2 要求分類 & Wiegers and Beatty; Sommerville \\
1.3 製品要求とプロジェクト要求 & Wiegers and Beatty \\
1.4 機能要求 & Wiegers and Beatty; Sommerville \\
1.5 非機能要求 & Wiegers and Beatty; Sommerville; ISO/IEC 25010 \\
1.6 技術制約 & Tockey, \textit{How to Engineer Software}; Yourdon \\
1.7 サービス品質制約 & ISO/IEC 25010; Tockey, \textit{Return on Software} \\
1.8 分類の理由 & Wiegers and Beatty; Tockey; McMenamin and Palmer \\
1.9 システム要求とソフトウェア要求 & INCOSE Handbook; ISO/IEC/IEEE 29148 \\
1.10 派生要求 & Tockey; INCOSE Handbook \\
1.11 要求活動 & Wiegers and Beatty; Sommerville \\
2. 要求獲得 & Wiegers and Beatty; Sommerville; Robertson and Robertson; LaPlante and Kassab \\
2.1 要求源泉 & Wiegers and Beatty; Robertson and Robertson \\
2.2 要求獲得技法 & Wiegers and Beatty; Wood and Silver; Gottesdiener; Brown and Katz \\
3. 要求分析 & Wiegers and Beatty; Sommerville; Tockey \\
3.1 要求の性質 & Wiegers and Beatty; Robertson and Robertson \\
3.2 サービス品質要求の経済性 & Tockey, \textit{Return on Software} \\
3.3 形式的分析 & Sommerville; Wing; Software Engineering Models and Methods \\
3.4 要求衝突への対応 & Robertson and Robertson; Fisher and Ury; Weiss and Lai \\
4. 要求仕様化 & Wiegers and Beatty; Sommerville; ISO/IEC/IEEE 29148 \\
4.1 非構造化自然言語 & Wiegers and Beatty; ISO/IEC/IEEE 29148 \\
4.2 構造化自然言語 & Wiegers and Beatty; Cockburn; Robertson and Robertson \\
4.3 受け入れ基準に基づく仕様化 & Sommerville; Smart; Software Testing KA \\
4.4 モデルに基づく仕様化 & Wiegers and Beatty; Sommerville; Wing; Tockey \\
4.5 要求の追加属性 & Wiegers and Beatty; Gilb \\
4.6 増分的仕様化と包括的仕様化 & Wiegers and Beatty; Robertson and Robertson \\
5. 要求妥当性確認 & Wiegers and Beatty; Sommerville; LaPlante and Kassab \\
5.1 要求レビュー & Wiegers and Beatty; Software Quality KA \\
5.2 シミュレーションと実行 & Tockey; Model-Based Specification references \\
5.3 プロトタイピング & Wiegers and Beatty; Sommerville \\
6. 要求管理活動 & Wiegers and Beatty; Sommerville; McConnell \\
6.1 要求整理・削減 & McConnell, \textit{Rapid Development} \\
6.2 要求変更管理 & Wiegers and Beatty; Sommerville; Software Configuration Management KA \\
6.3 範囲調整 & McConnell; Software Engineering Management KA \\
7. 実務上の考慮 & Wiegers and Beatty; Sommerville; Tockey \\
7.1 反復性 & Sommerville; Robertson and Robertson \\
7.2 優先順位付け & Wiegers and Beatty \\
7.3 追跡可能性 & Wiegers and Beatty; Gotel and Finkelstein \\
7.4 安定性と変動性 & Sommerville; Tockey \\
7.5 要求測定 & Wiegers and Beatty; Software Engineering Process KA \\
7.6 プロセス品質と改善 & Wiegers and Beatty; Software Quality KA \\
8. ソフトウェア要求ツール & Wiegers and Beatty; Sommerville \\
8.1 要求管理ツール & Wiegers and Beatty \\
8.2 要求モデリングツール & Wiegers and Beatty; Sommerville; Formal Methods references \\
8.3 機能テストケース生成ツール & Software Testing KA; Tockey \\
\bottomrule
\end{longtable}
\normalsize

\section{発展的な読み物}

SWEBOK第01章で紹介される発展的な読み物を、日本語で読む目的に合わせて整理する。

\begin{tabularx}{0.97\linewidth}{@{}p{.31\linewidth}X@{}}
\toprule
文献 & 読むと得られること \\
\midrule
\textit{BABOK Guide} & ビジネス分析全体の知識体系。ソフトウェア要求だけでなく、ビジネス課題、利害関係者、価値の分析を広く学べる。 \\
LaPlante and Kassab & 要求工学を体系的に学ぶための包括的な教科書。 \\
Robertson and Robertson & 要求プロセス、利害関係者、要求発見、要求記述を実務的に学べる。 \\
Gilb & 要求を測定可能にし、価値と品質を重視して記述する考え方を学べる。 \\
Wiegers, \textit{Software Development Pearls} & 長年の実務経験から得られた、要求を含むソフトウェア開発の実践的教訓を学べる。 \\
Fisher and Ury & 要求衝突の解決に役立つ交渉の基本を学べる。 \\
Ahmad & 電子的コミュニケーションを用いた面接で、暗黙知を引き出す方法に関する研究。 \\
\bottomrule
\end{tabularx}

\section{参考文献}

\footnotesize
\begin{enumerate}[leftmargin=2.2em,itemsep=.15ex]
\item K. E. Wiegers and J. Beatty, \textit{Software Requirements}, 3rd ed., Microsoft Press, 2013.
\item I. Sommerville, \textit{Software Engineering}, 10th ed., Addison-Wesley, 2016.
\item S. Tockey, \textit{Return on Software: Maximizing the Return on Your Software Investment}, Addison-Wesley, 2005.
\item J. M. Wing, ``A Specifier's Introduction to Formal Methods,'' \textit{Computer}, 1990.
\item P. Laplante and M. Kassab, \textit{Requirements Engineering for Software and Systems}, 4th ed., CRC Press, 2022.
\item S. Robertson and J. Robertson, \textit{Mastering the Requirements Process: Getting Requirements Right}, Addison-Wesley, 2013.
\item T. Gilb, \textit{Competitive Engineering}, Elsevier Butterworth-Heinemann, 2005.
\item E. Yourdon, \textit{Modern Structured Analysis}, Prentice-Hall, 1989.
\item S. Tockey, \textit{How to Engineer Software}, Wiley, 2019.
\item S. Ambler, \textit{Agile Modeling}, Wiley, 2002.
\item A. Cockburn, \textit{Writing Effective Use Cases}, Addison-Wesley, 2000.
\item L. Constantine and L. Lockwood, \textit{Software for Use}, Addison-Wesley, 2000.
\item J. Wood and D. Silver, \textit{Joint Application Development}, Wiley, 1995.
\item E. Gottesdiener, \textit{Requirements by Collaboration}, Addison-Wesley, 2002.
\item J. Terninko, \textit{Step by Step QFD}, 2nd ed., CRC Press, 1997.
\item G. Salvendy, \textit{Handbook of Human Factors}, 4th ed., Wiley, 2012.
\item T. Brown and B. Katz, \textit{Change by Design}, Harper Collins, 2019.
\item S. McMenamin and J. Palmer, \textit{Essential Systems Analysis}, Yourdon Press, 1984.
\item J. Smart, \textit{BDD in Action}, Manning Publications, 2015.
\item D. Weiss and C. Lai, \textit{Software Product-Line Engineering}, Addison-Wesley, 1999.
\item K. Wiegers, \textit{Software Development Pearls}, Addison-Wesley Professional, 2021.
\item S. McConnell, \textit{Rapid Development}, Microsoft Press, 1996.
\item O. Gotel and C. W. Finkelstein, ``An Analysis of the Requirements Traceability Problem,'' 1994.
\item INCOSE, \textit{Systems Engineering Handbook}, International Council on Systems Engineering, 2012.
\item R. Fisher and W. Ury, \textit{Getting to Yes}, 3rd ed., Penguin, 2011.
\item ISO/IEC/IEEE 29148, \textit{Systems and software engineering -- Life cycle processes -- Requirements engineering}, 2018.
\item ISO/IEC 25010, \textit{Systems and software Quality Requirements and Evaluation (SQuaRE) -- System and software quality models}, 2011.
\item ISO/IEC/IEEE 24765:2017, \textit{Systems and Software Engineering -- Vocabulary}, 2nd ed., 2017.
\item N. Ahmad, \textit{Effects of Electronic Communication on the Elicitation of Tacit Knowledge in Interview Techniques for Small Software Developments}, doctoral thesis, University of Huddersfield, 2021.
\item IIBA, \textit{A Guide to the Business Analysis Body of Knowledge (BABOK Guide) v3}, 2015.
\end{enumerate}
\normalsize

\section{画像TODO一覧}

本文に配置した画像TODOを再掲する。実際に図を作る場合は、各TODOのプロンプトをそのまま画像生成に使うと、A4資料に合う図を作りやすい。

\begin{enumerate}
\item 第01章の全体地図。
\item 要求分類の階層図。
\item 完全技術フィルタの判定フロー。
\item システム要求からソフトウェア要求への割当図。
\item 要求開発と要求管理の活動図。
\item 要求の源泉と獲得技法の対応図。
\item サービス品質制約の価値・費用曲線。
\item 要求仕様化技法の比較図。
\item 要求妥当性確認の三手法比較図。
\item 要求変更管理の流れ。
\item 要求追跡の連鎖図。
\item 要求ツールの分類図。
\end{enumerate}

\section{章末確認問題}

\begin{enumerate}
\item 「利用者は注文番号で注文履歴を検索できる」は、どの分類の要求か。
\item 「検索結果は通常時2秒以内に表示される」は、機能要求か非機能要求か。理由も述べよ。
\item 「実装言語はPythonとする」は、技術制約かサービス品質制約か。
\item 要求獲得と要求分析の違いを、自分の言葉で説明せよ。
\item 要求仕様化において、自然言語、受け入れ基準、モデルの違いを説明せよ。
\item 要求妥当性確認で、レビュー、シミュレーション、プロトタイプはどのように使い分けるか。
\item 要求変更を受け入れるとき、なぜ費用・納期・スコープへの影響を確認する必要があるか。
\item 追跡可能性が、変更時の影響分析に役立つ理由を説明せよ。
\end{enumerate}

\MiniBox{解答の方向性}{1は機能要求。2はサービス品質制約であり非機能要求。3は技術制約。4以降は、本文の「合意を作る活動」「候補要求を表面化させる」「意味・矛盾・実現可能性を調べる」「合意を維持する」といった語を使って説明できる。}

\end{document}
